\section{Multi-Specialty Agent Decomposition Engine}
\label{sec:multi_agent_engine}

The core cognitive engine of MedRAGAgents relies on an ensemble of specialized domain agents working in concert to analyze clinical vignettes. Rather than forcing a single foundation model prompt to perform multi-domain reasoning simultaneously, the workload is distributed across specialized reasoning components.

\subsection{Domain Classification Agent (\texttt{DomainAgent})}

The \texttt{DomainAgent} acts as the front-end classifier for clinical vignettes. Its objective is to analyze the raw clinical text and extract relevant medical subspecialties for both the vignette body and candidate options.

Figure~\ref{fig:domain_routing} illustrates the domain classification workflow.

\begin{figure}[htbp]
\centering
\resizebox{\linewidth}{!}{%
\begin{tikzpicture}[node distance=1.0cm and 1.5cm, auto]
  \node [monobox] (vignette) {Clinical Vignette +\\Options Payload};
  \node [monobox, right=of vignette] (domain_agent) {\texttt{DomainAgent}\\Classifier};
  
  \node [monobox, above right=0.6cm and 1.2cm of domain_agent] (q_domains) {Question Domains\\(Max 5 Specialties)};
  \node [monobox, below right=0.6cm and 1.2cm of domain_agent] (o_domains) {Option Domains\\(Max 2 Specialties)};
  
  \node [monobox, right=1.2cm of q_domains] (q_store) {Store in \texttt{ShortTermMemory}\\(\texttt{"q\_domains"})};
  \node [monobox, right=1.2cm of o_domains] (o_store) {Store in \texttt{ShortTermMemory}\\(\texttt{"o\_domains"})};

  \draw [monoline] (vignette) -- (domain_agent);
  \draw [monoline] (domain_agent.east) |- (q_domains.west);
  \draw [monoline] (domain_agent.east) |- (o_domains.west);
  \draw [monoline] (q_domains) -- (q_store);
  \draw [monoline] (o_domains) -- (o_store);
\end{tikzpicture}%
}
\caption{Domain Classification Routing Workflow.}
\label{fig:domain_routing}
\end{figure}

The classification process executes in two distinct operations:
\begin{enumerate}
    \item \textbf{Question Domain Classification} (\texttt{classify\_question\_domains}): Maps the clinical vignette to up to five primary medical domains (e.g., Cardiology, Nephrology, Medical Microbiology, Pharmacology, Pathology).
    \item \textbf{Option Domain Classification} (\texttt{classify\_option\_domains}): Analyzes candidate choices against the vignette context to identify up to two secondary specialties relevant to option differentiation.
\end{enumerate}

\subsection{Specialized Domain Analysis Agents (\texttt{AnalysisAgent})}

Once subspecialties are identified, the \texttt{AnalysisAgent} executes isolated specialty evaluations. This stage prevents cross-domain context contamination by restricting each prompt invocation to a single specialty perspective.

Figure~\ref{fig:parallel_dag} displays the execution Directed Acyclic Graph (DAG) for parallel specialty analyses.

\begin{figure}[htbp]
\centering
\resizebox{\linewidth}{!}{%
\begin{tikzpicture}[node distance=1.0cm and 1.5cm, auto]
  \node [monobox] (active_domains) {Active Domain Set\\$\mathcal{D}_{\text{active}}$};
  
  \node [monobox, right=1.5cm of active_domains, yshift=1.5cm] (agent1) {Pathology Expert\\Analysis};
  \node [monobox, right=1.5cm of active_domains, yshift=0.5cm] (agent2) {Pharmacology Expert\\Analysis};
  \node [monobox, right=1.5cm of active_domains, yshift=-0.5cm] (agent3) {Neurology Expert\\Analysis};
  \node [monobox, right=1.5cm of active_domains, yshift=-1.5cm] (agent4) {Option-Specific\\Domain Analysis};
  
  \node [monobox, right=2.0cm of agent2, yshift=-0.5cm] (aggregator) {Analysis Dictionary\\Aggregation};

  \draw [monoline] (active_domains.east) |- (agent1.west);
  \draw [monoline] (active_domains.east) |- (agent2.west);
  \draw [monoline] (active_domains.east) |- (agent3.west);
  \draw [monoline] (active_domains.east) |- (agent4.west);
  
  \draw [monoline] (agent1.east) -| (aggregator.north);
  \draw [monoline] (agent2.east) -- (aggregator.west);
  \draw [monoline] (agent3.east) -- (aggregator.west);
  \draw [monoline] (agent4.east) -| (aggregator.south);
\end{tikzpicture}%
}
\caption{Multi-Agent Parallel Specialty Analysis Execution DAG.}
\label{fig:parallel_dag}
\end{figure}

The analysis phase operates in two stages:
\begin{itemize}
    \item \textbf{Question Specialty Analysis}: Each question domain specialist evaluates the vignette independently, extracting pertinent clinical findings, lab abnormalities, and pathophysiological mechanisms.
    \item \textbf{Option Specialty Analysis}: Option specialists evaluate each candidate option against the vignette and the prior question analyses, assessing diagnostic plausibility and contraindications.
\end{itemize}

\subsection{Diagnostic Synthesis Engine (\texttt{SynthesisAgent})}

The \texttt{SynthesisAgent} aggregates outputs from all specialty analysis agents along with retrieved textbook evidence snippets, producing a unified diagnostic report.

Figure~\ref{fig:synthesis_flow} illustrates the evidence integration flow.

\begin{figure}[htbp]
\centering
\resizebox{\linewidth}{!}{%
\begin{tikzpicture}[node distance=1.0cm and 1.5cm, auto]
  \node [monobox] (q_ana) {Question Specialty\\Analyses};
  \node [monobox, below=0.6cm of q_ana] (o_ana) {Option Specialty\\Analyses};
  \node [monobox, below=0.6cm of o_ana] (rag_snip) {MedCPT Retrieved\\Textbook Snippets};
  
  \node [monobox, right=1.5cm of o_ana] (synth) {\texttt{SynthesisAgent}\\Engine};
  \node [monobox, right=1.5cm of synth] (report) {Unified Diagnostic\\Synthesis Report};

  \draw [monoline] (q_ana.east) -| (synth.north);
  \draw [monoline] (o_ana.east) -- (synth.west);
  \draw [monoline] (rag_snip.east) -| (synth.south);
  \draw [monoline] (synth.east) -- (report.west);
\end{tikzpicture}%
}
\caption{Diagnostic Synthesis and Evidence Fusion Architecture.}
\label{fig:synthesis_flow}
\end{figure}

The synthesis report presents a complete diagnostic breakdown, including:
\begin{enumerate}
    \item Summary of key clinical findings and lab data.
    \item Differential diagnosis ranking across candidate options.
    \item Explicit rejection rationales for incorrect distractors.
    \item Definitive option selection based on combined clinical evidence.
\end{enumerate}

\subsection{Consensus Verification and Voting Protocol (\texttt{VerifierAgent})}

The \texttt{VerifierAgent} manages an iterative peer-verification consensus loop. Rather than accepting the synthesis report without validation, specialty agents review the report, vote on its diagnostic validity, and request revisions if discrepancies exist.

Figure~\ref{fig:verifier_fsm} depicts the finite state machine governing consensus verification.

\begin{figure}[htbp]
\centering
\resizebox{\linewidth}{!}{%
\begin{tikzpicture}[node distance=1.0cm and 1.5cm, auto]
  \node [monobox] (start) {Synthesis Report\\Generated};
  \node [monobox, right=of start] (vote) {Specialty Voting\\(\texttt{\_vote()})};
  \node [monodiamond, right=of vote] (check) {Unanimous\\Consensus?};
  
  \node [monobox, above=1.2cm of check] (advice) {Collect Revision\\Advice (\texttt{\_get\_advice()})};
  \node [monobox, left=of advice] (revise) {Revise Report\\(\texttt{\_revise()})};
  
  \node [monobox, right=1.5cm of check] (final) {Extract Final Option\\(\texttt{\_derive\_final\_answer()})};

  \draw [monoline] (start) -- (vote);
  \draw [monoline] (vote) -- (check);
  \draw [monoline] (check) -- node [above] {YES} (final);
  
  \draw [monoline] (check) -- node [right] {NO ($t < 3$)} (advice);
  \draw [monoline] (advice) -- (revise);
  \draw [monoline] (revise.south) -- (vote.north);

  \draw [monoline] (check.south) -- ++(0,-1.0) -| (final.south) node[pos=0.25, below] {Max Rounds ($t=3$)};
\end{tikzpicture}%
}
\caption{Consensus Verification Finite State Machine Loop.}
\label{fig:verifier_fsm}
\end{figure}

The verification protocol executes as follows:
\begin{enumerate}
    \item \textbf{Voting Pass} (\texttt{\_vote}): Each active specialty domain node reviews the current synthesis report and submits a binary vote (\texttt{YES} / \texttt{NO}).
    \item \textbf{Consensus Check}: If all domain nodes vote \texttt{YES}, the consensus loop terminates successfully.
    \item \textbf{Advice Generation} (\texttt{\_get\_advice}): If any domain node votes \texttt{NO}, dissenting domain nodes provide targeted revision recommendations identifying factual errors or logical oversights.
    \item \textbf{Report Revision} (\texttt{\_revise}): The synthesis report is updated based on dissenting feedback, and the verification loop increments to round $t+1$.
    \item \textbf{Termination and Derivation} (\texttt{\_derive\_final\_answer}): The loop terminates upon reaching unanimous consensus or completing $t=3$ rounds. The final option letter (A--E) is extracted via structured regex parsing.
\end{enumerate}
